\bookchapter{第五章：スマートフォンをもっと理解させる——基本設定}

\chapterintro{数分かけて設定するだけで、体験は大幅に向上する。これらの設定はスマートフォンに「チューニング」を行うようなもので、あなたの使用習慣に合わせることができる。}

\booksection{一、サウンドとバイブレーション一、サウンドとバイブレーション}

：聞こえて、かつ恥ずかしくないものを選ぶ。

\begin{itemize}
  \item \textbf{バイブレーションモード}：マナー時にバイブで通知。
  \item \textbf{おやすみモード}：例外を設定可能。連絡先からの着信を許可するなど。
  \item \textbf{音量分割}：メディア音量、着信音量、アラーム音量を個別に調整可能。
  \item \textbf{二、表示と明るさ}自動明るさ
\end{itemize}

\booksection{二、表示と明るさ二、表示と明るさ}

\begin{itemize}
  \item \textbf{アイケア/ナイトシフトモード}：ブルーライトをカットし、夜間の閲覧を快適に。
  \item \textbf{ダークモード}：省電力（OLED画面で特に有効）、かつ夜間に眩しくない。多くのアプリもダークモードをサポート。
  \item \textbf{画面リフレッシュレート}：一部のAndroidスマートフォンは調整可能。高いほど（120Hzなど）スクロールが滑らか。
  \item \textbf{自動ロック}：15秒または30秒に設定。安全で省電力。
  \item \textbf{フォントサイズ}：字が小さいと感じる？少し大きくしよう。
  \item \textbf{三、壁紙、テーマとホーム画面}素敵な壁紙（静的または動的）に変更するだけで、気分も変わる。Androidブランドのスマートフォンにはテーマストアがあり、アイコンやスタイルを丸ごと変更可能。
\end{itemize}

\booksection{三、壁紙、テーマとホーム画面三、壁紙、テーマとホーム画面}

：天気、スケジュール、音楽コントロールをホーム画面に配置し、ワンタップでアクセス。iOSとAndroidの両方がサポート。ホーム画面の空白部分を長押しして追加。

\textbf{四、効率設定：スマートフォンを「外付け脳」にする}このセクションは「一度設定したら戻れない」機能で、どれも日常効率を大幅に向上させる。

\booksection{四、効率設定：スマートフォンを「外付け脳」にする四、効率設定：スマートフォンを「外付け脳」にする}

核心となる考え

\booksubsection{集中 / おやすみモード：通知に奪われた注意力を取り戻す}

\begin{itemize}
  \item \textbf{使用のアドバイス}：
  \item \textbf{· 仕事モード：「よく連絡する」連絡先と必要なアプリのみ許可；}：\newline · 運動 / 読書 / 会議：現在の活動に応じてカスタマイズ。\newline 応用テクニック\newline ：自動化 + 位置 / 時間トリガーでモードを自動切替。例：iOSショートカットで「会社に到着したら自動的に仕事モードに切り替え」を設定可能。
  \item \textbf{2. 通知管理：赤いバッジの「しつこい連呼」を止める}マーケティング通知をオフ
\end{itemize}

\booksubsection{通知管理：赤いバッジの「しつこい連呼」を止める}

\begin{itemize}
  \item \textbf{重要度で並べ替え}：人間系の通信（電話、SMS、インスタントメッセージ）をトップに固定；ニュース、ショッピング、ソーシャルプラットフォーム系は「サイレントプッシュ」に設定可能。
  \item \textbf{グループ化通知}：同種の通知をまとめて折りたたみ、一度のロック解除で数十件のメッセージに埋もれるのを防ぐ。
  \item \textbf{3. ショートカットジェスチャー：操作を一歩リード}スクリーンショット
\end{itemize}

\booksubsection{ショートカットジェスチャー：操作を一歩リード}

\begin{itemize}
  \item \textbf{長尺スクリーンショット / スクロールキャプチャ}：通常のスクリーンショット後に「長尺スクリーンショット」オプションをタップ。
  \item \textbf{ワンタップQRスキャン}：コントロールセンター / クイック設定に「スキャン」アイコンを追加。
  \item \textbf{背面タップ（iPhone）}：[設定] > [アクセシビリティ] > [タッチ] > [背面タップ]で「スクリーンショット」や「ロック」などのアクションを割り当て。
  \item \textbf{音声起動}：Siri / 小愛同学 / 小芸 / 小布、一声かければ天気確認やアラーム設定が可能。
  \item \textbf{4. 入力メソッド：最も過小評価されている効率化ツール}適切な入力メソッドの選択
\end{itemize}

\booksubsection{入力メソッド：最も過小評価されている効率化ツール}

\begin{itemize}
  \item \textbf{· Sogou / Baidu：辞書が大きく、流行語の更新が速い；}：\newline · Gboard（Google）：シンプル、クロスプラットフォーム同期が強力、フリック入力をサポート；\newline · iOS標準：プライバシー保護が最も良く、システムと深く統合。\newline 必須のスイッチ\newline ：スマート補正、予測変換、クラウド辞書同期、上スワイプで数字·記号入力。
  \item \textbf{プライバシーの注意}：「文章予測」や「個別化推奨」機能は慎重に使用。入力内容がアップロードされる可能性がある。
  \item \textbf{5. ホーム画面とマルチタスク：ワークフローを滑らかに}アプリ分割表示 / ピクチャーインピクチャー
\end{itemize}

\booksubsection{ホーム画面とマルチタスク：ワークフローを滑らかに}

\begin{itemize}
  \item \textbf{小ウィンドウ / フローティングウィンドウ}：WeChat / TikTokは「小ウィンドウモード」をサポートし、返信後にすぐ閉じられる。
  \item \textbf{常用アプリを固定}：Dockバー（画面下部）には最も高頻度の4つのみ配置し、検索時間を削減。
  \item \textbf{「初日勤務」チェックリスト}新規スマートフォンの取得や大型アップデート後、10分かけて以下の設定を確認することを推奨：
\end{itemize}

\begin{tcbinfo}{(1) ロック画面パスワード + 生体認証 (2) 通知グループ化 + マーケティングプッシュをオフ (3) フォントサイズ / リフレッシュレート / ダークモード}
(4) デフォルトの入力メソッド / ブラウザ / メール (5) 緊急連絡先 / メディカルID (6) バックアップ設定\newline 五、アクセシビリティ機能：テクノロジーが善である力\newline 「アクセシビリティ機能」という名前に恐れる必要はない。これらの機能は誰でも使える。元々は障害者のために設計されたが、一般人にも役立つ——一時的に手を怪我した、赤ちゃんを抱えている、地下鉄でイヤホンが壊れて聞こえない、老眼鏡で小さな文字を見る……。どのような状況でも助けになる。結局のところ、アクセシビリティデザインは

\end{tcbinfo}

\booksection{五、アクセシビリティ機能：テクノロジーが善である力五、アクセシビリティ機能：テクノロジーが善である力}

理念を体現している：すべての人が平等にデジタルライフを享受できるようにする。\textbf{1. 視覚障害者の「目」——スクリーンリーダー完全ガイド}スクリーンリーダーは視覚障害者がスマートフォンを使用するためのコアツール。画面上の

\booksubsection{視覚障害者の「目」——スクリーンリーダー完全ガイド}

、ユーザーはジェスチャーと音声でスマートフォンを操作する。以下ではiOS「VoiceOver」とAndroid「TalkBack」の使用方法をそれぞれ紹介する。\textbf{iOS VoiceOver}有効化の手順

\begin{tcbinfo}{：[設定] > [アクセシビリティ] > [VoiceOver]、スイッチをオン。サイドボタン（またはホームボタン）を3回クリックして素早く切り替えも可能。}
\textbf{基本ジェスチャー}：

\textbf{指1本で左右スワイプ：フォーカス要素を切り替え；}：

\begin{enumerate}
  \item 指1本でダブルタップ：現在のフォーカス要素をアクティブ化（「タップ」に相当）；
  \item 指3本でスワイプ：ページを切り替え；
  \item 指1本でトリプルタップ：拡大または選択（通常のトリプルタップとは異なる）。
  \item 個別設定
  \item ：[VoiceOver] > [音声]で、
\end{enumerate}

\textbf{読み上げ速度}（デフォルト約50\%、70-80\%に調整すると聴感がよりスムーズ）、\textbf{ピッチ}、異なる音声パッケージの切り替えが可能。iOSには「拡張音声」と「音声コントローラ」が内蔵され、前者は呼吸音や間を減らし、より自然な聴感を実現。\textbf{音声エンジン}有効化の手順

\end{tcbinfo}

\begin{tcbinfo}{：[設定] > [アクセシビリティ] > [TalkBack]、スイッチをオン。または音量上下キーを3秒間長押しで素早く起動（一部ブランドがサポート）。}
\textbf{基本ジェスチャー}：

\textbf{指1本で左右スワイプ：フォーカス要素を切り替え；}：

\begin{enumerate}
  \item 指2本で上下/左右スワイプ：スクロール；
  \item 指3本でシングルタップ：読み上げモード（連続読み上げ）に移行；
  \item 指4本でシングルタップ：読み上げを一時停止/再開。
  \item 個別設定
  \item ：[TalkBack] > [テキスト読み上げ設定]で、
\end{enumerate}

\textbf{読み上げ速度}、\textbf{ピッチ}、\textbf{音声エンジン}の切り替えが可能（Googleテキスト読み上げエンジンのインストールを推奨、より自然な読み上げ）。\textbf{2. 小さな字を大きく、ぼやけを鮮明に——視覚調整}拡大ジェスチャー

\end{tcbinfo}

\booksubsection{小さな字を大きく、ぼやけを鮮明に——視覚調整}

\booksubsection{拡大ジェスチャー}

\begin{itemize}
  \item \textbf{iOS}で拡大。デフォルトは全画面ズーム、ウィンドウズーム（部分拡大ウィンドウ）も選択可能。\textbf{2本指のピンチで倍率調整（1.5x〜15x）}、3本指でダブルタップすると現在のズーム倍率を固定。\textbf{ピンチイン·ピンチアウトで倍率を調整（1.5倍〜15倍）}：[設定] > [アクセシビリティ] > [拡大]、3本指でタップして拡大、もう一度3本指でタップして終了；または
  \item \textbf{Android}して一時拡大を有効化（最大8倍）。\textbf{ハイコントラストテキスト / 太字テキスト}iOS
\end{itemize}

\booksubsection{ハイコントラストテキスト / 太字テキスト}

\begin{itemize}
  \item \textbf{iOS}：[設定] > [表示] > [フォントサイズ]、「小、標準、大、特大」の4段階を調整可能。Huawei/Xiaomiは「太字」と「ハイコントラストテキスト」スイッチを追加提供。
  \item \textbf{Android}iOS
\end{itemize}

\booksubsection{カラーフィルター（色覚バリアフリー）}

\begin{itemize}
  \item \textbf{iOS}、および「グレースケール」と「色を反転」オプションを提供。まず色覚テスト（システム内蔵）を行ってから対応するフィルターを選択することを推奨。\textbf{赤／緑／青／黄の4つの色覚サポートモード}：[設定] > [アクセシビリティ] > [色補正]で、
  \item \textbf{Android}+ 「色を反転」スイッチを含む。\textbf{3. AI画像記述：AIをあなたの「目」にする}スクリーンリーダーが画像に遭遇した場合、従来はスキップするか「画像」という2文字だけを読み上げていた。しかし今、AIは
\end{itemize}

\booksubsection{AI画像記述：AIをあなたの「目」にする}

でき、認証コードの識別、画像の読み取り、メニューの識別に非常に役立つ。\textbf{各システムのAI画像記述サポート状況}iOS 15+

\begin{tcbinfo}{：[設定] > [アクセシビリティ] > [VoiceOver] > [音声記述]で「画像記述」をオン。システムは}
\begin{itemize}
  \item \textbf{ローカルAIモデル}を呼び出して画像のテキスト記述を自動生成（iPhone 12以降が必要）。\textbf{新版TalkBackアプリ}：Google TalkBackは比較的新しいバージョンで
  \item \textbf{画像Q\& A}機能を統合。「この画像に何がある？」と質問すると、AIがリアルタイムに回答（Androidシステムバージョンに厳密に依存せず、TalkBackアプリを更新するだけで体験可能）。\textbf{サードパーティアプリ}：「Seeing AI」（iOS）、「Lookout」（Android）、「小芸幫幫」（Huawei）はいずれも優れたAI視覚アシスタントで、物品認識、書類読み取り、顔の表情認識が可能。
  \item \textbf{使用シーン}街撮り写真に向かって「この画像には何人がいますか？彼らは何をしていますか？」と聞いてみよう。またはデリバリーアプリを開き、AIにメニュー画像の料理名を読み上げさせよう——
\end{itemize}

\end{tcbinfo}

\begin{tcbpractice}{これこそが真の「テクノロジーが生活をより平等にする」}
。\textbf{4. 聴覚障害者の「耳」——リアルタイム字幕完全ガイド}。

\end{tcbpractice}

\booksubsection{聴覚障害者の「耳」——リアルタイム字幕完全ガイド}

機能で、聴覚障害者が騒がしい環境、図書館、外国語シーンで特に役立つ。\textbf{リアルタイム字幕の有効化と最適化}iOS

\begin{tcbinfo}{：[設定] > [アクセシビリティ] > [リアルタイム字幕]、スイッチをオン。システムはデバイスのすべての音声（通話、動画、ポッドキャスト、FaceTime、ボイスメモ）を検出し、自動的に画面上部に字幕を表示。}
\textbf{iOS}：

\textbf{音量キー}を押す ポップアップボックスの「リアルタイム字幕」アイコンをタップ；または[設定] > [アクセシビリティ] > [リアルタイム字幕]でオン。\textbf{精度向上のヒント}：

\textbf{できるだけ}：

\begin{enumerate}
  \item で使用し、エアコンや機械の轟音から離れる；\textbf{話者が}はっきり発音し、適切な速度
  \item の時に認識率が最も高い；\textbf{複数人の対話は精度を下げる、「1人だけ表示」モードを使用可能；}中国語普通話、広東語、英語などの主要言語をサポート、
  \item 少数言語の認識率は低い
  \item ；\textbf{話の途中にBGMや人の声のノイズがあると、認識が途切れ途切れになる。}5. 通話シーンのデュアルチャネル：音声テキスト変換 + テキスト音声変換
  \item ユーザーが
\end{enumerate}

\end{tcbinfo}

\booksubsection{通話シーンのデュアルチャネル：音声テキスト変換 + テキスト音声変換}

場合（突発性失声、喉の手術回復期、重度の吃音など）、スマートフォンは「伝言役」になる。\textbf{通話リアルタイム翻訳（TTY / RTT）}iOS

\begin{tcbinfo}{：[設定] > [アクセシビリティ] > [RTT/TTY]、有効化すると通話中に画面にキーボードが表示され、相手の音声がリアルタイムにテキスト化され、}
\textbf{iOS}。\textbf{入力した文字も相手に読み上げられます}：[設定] > [アクセシビリティ] > [リアルタイム文字起こし] + [音声通知]。「リアルタイム文字起こし」は相手の発話をリアルタイムにテキスト表示；「音声通知」は音声が不便な時にあなたの入力を音声で再生。

\textbf{Android}：

\textbf{USB PD 3.1}：

\begin{itemize}
  \item \textbf{TTY（テレタイプライター）}：専用TTYデバイスが必要、RTTに代替されつつある；
  \item \textbf{WeChat/QQの}音声テキスト変換機能
  \item も代替案：音声メッセージを長押し 「テキスト変換」。\textbf{6. すべての人向けの「小さくても実用的な」機能}背面タップ
\end{itemize}

\end{tcbinfo}

\booksubsection{すべての人向けの「小さくても実用的な」機能}

\begin{itemize}
  \item \textbf{ダブルタップ/トリプルタップ}でスクリーンショット、ロック、コントロールセンターなどのアクションを割り当て可能。両手で荷物を持っている時に特に便利。\textbf{アシストタッチ / フローティングボール}：iOS [アクセシビリティ] > [タッチ] > [アシストタッチ]、Android [アクセシビリティ] > [フローティングナビゲーション]。仮想ボタンでジェスチャーや物理ボタンを代替、ボタン不調や片手操作時に助けに。
  \item \textbf{スイッチコントロール}：外部スイッチ、シングルタップ、頭部の動作でスマートフォンを操作、
  \item \textbf{肢体障害者向けに設計}。iOS [アクセシビリティ] > [スイッチコントロール]、Android [アクセシビリティ] > [スイッチコントロール]。\textbf{サウンド認識}（iOS）：[アクセシビリティ] > [サウンド認識]、有効化するとスマートフォンがドアベル、赤ちゃんの泣き声、火災警報、サイレンなどの音を「聞き」、
  \item \textbf{認識するとバイブレーション+画面通知}。聴覚障害者や睡眠中に非常に有用。\textbf{LEDフラッシュ通知}：iOS [アクセシビリティ] > [オーディオ/視覚] > [LEDフラッシュで通知]、Android [設定] > [通知] > [フラッシュライト通知]。着信や通知時にフラッシュライトが点滅、
  \item \textbf{騒がしい環境でも音なしで感知可能}。\textbf{モノラルオーディオ}。
  \item \textbf{Appleのアクセシビリティデザイン哲学}Appleはアクセシビリティ分野に多大な投資を行い、毎年アクセシビリティ機能の更新と追加を行っている。Apple公式サイトには専用のアクセシビリティページがあり、各機能の使用方法を詳しく紹介。iOSのアクセシビリティ機能は通常Androidより完璧で、体験もよりスムーズ。これが多くの障害者がAppleを選ぶ重要な理由だ。
\end{itemize}

\begin{tcbinfo}{30秒で体験してみよう}
空き時間を見つけて、以下の体験をしよう：

\end{tcbinfo}

\begin{tcbpractice}{VoiceOver / TalkBack}
をオンにし、目に頼らず3本指でスワイプしてWeChatを閲覧してみる；

\begin{enumerate}
  \item リアルタイム字幕\textbf{をオンにし、スマートフォンに向かって話して認識精度を確認；}3本指タップ
  \item でテキストを200\%に拡大し、視覚障害者の日常を体験；\textbf{すべての機能をオフにしてから、「もし私が永遠にこれらの機能を必要とするなら、どう設計してほしいか」を考える。}第六章：デジタルライフを守る——セキュリティとプライバシー
  \item スマートフォンのセキュリティ＝プライバシーセキュリティ＋資産セキュリティ。データが資産となる時代において、スマートフォンを守ることは自分自身を守ること。\textbf{一、ロック画面：第一の物理的防線}ロック画面はスマートフォンの第一の防線であり、最も見落とされやすい。
  \item セキュリティレベルの順位：
\end{enumerate}

\end{tcbpractice}
